\chapter{第六单元}

\Needspace{12\baselineskip}

\section{第六单元素养评价（A）}

\Needspace{8\baselineskip}

\pdfbookmark[2]{原书第49页}{orig-49}

\subsection*{原书第49页}

% p049_s01
\BookQuestion{一、1}
中国国家博物馆遵循“开门办馆”的理念，分期分批对社会公布馆藏藏品目录，第一期共公布藏品约30.64万件。横线上的小数读作（\hspace{1.8em}），这个小数的计数单位是（\hspace{1.8em}），它有（\hspace{1.8em}）个这样的计数单位。\par
\begin{BookAnswer}
依次填“三十点六四、\(0.01\)、\(3064\)”。整数部分30按整数读法读“三十”，小数部分按数位逐个读“六、四”。30.64写到百分位，计数单位是百分之一，即0.01。3个十相当于3000个0.01，6个十分之一相当于60个0.01，再加4个0.01，共\(3000+60+4=3064\)个。
\end{BookAnswer}

% p049_s02
\BookQuestion{一、2}
一个小数由8个0.1和7个0.01组成，这个小数是（\hspace{1.8em}），再加上（\hspace{1.8em}）个\(\frac1{100}\)就是1。\par
\begin{BookAnswer}
依次填\(0.87,13\)。\(8\)个\(0.1\)是\(0.8\)，\allowbreak{}\(7\)个\(0.01\)是\(0.07\)，合为\(0.87\)。把1写成\(1.00\)，还差\(1.00-0.87=0.13\)，即13个\(\frac1{100}\)。
\end{BookAnswer}

% p049_s03
\BookQuestion{一、3}
把168000改写成用“万”作单位的数是（\hspace{1.8em}）万，保留整数约是（\hspace{1.8em}）万；把995000000改写成以“亿”为单位的数是（\hspace{1.8em}）亿，精确到十分位约是（\hspace{1.8em}）亿。\par
\begin{BookAnswer}
依次填\(16.8,17,9.95,10.0\)。先准确改写，每10000合成1万、每100000000合成1亿：\[\begin{aligned}168000&=(168000\div10000)\text{万}=16.8\text{万}\\995000000&=(995000000\div100000000)\text{亿}=9.95\text{亿}\end{aligned}\]再求近似数。16.8保留整数看十分位8，向个位进1；9.95保留一位小数看百分位5，连续进位。\[\begin{aligned}16.8\text{万}&\approx17\text{万}\\9.95\text{亿}&\approx10.0\text{亿}\end{aligned}\]准确改写用等号，求近似数用约等号。10.0末尾的0必须保留，表示保留一位小数；原空后已有单位，只填数。
\end{BookAnswer}

% p049_s04
\BookQuestion{一、4}
3分米＝（\hspace{1.8em}）米；5.06吨＝（\hspace{1.8em}）吨（\hspace{1.8em}）千克；32克＝（\hspace{1.8em}）千克；5平方米3平方分米＝（\hspace{1.8em}）平方米。\par
\begin{BookAnswer}
依次填\(0.3;5,60;0.032;5.03\)。分米化米除以10，吨化千克乘1000，克化千克除以1000，平方分米化平方米除以100。每项换算另行写清单位：\[\begin{aligned}3\text{分米}&=(3\div10)\text{米}=0.3\text{米}\\0.06\text{吨}&=(0.06\times1000)\text{千克}=60\text{千克}\\5.06\text{吨}&=5\text{吨}60\text{千克}\\32\text{克}&=(32\div1000)\text{千克}=0.032\text{千克}\\3\text{平方分米}&=(3\div100)\text{平方米}=0.03\text{平方米}\\5\text{平方米}3\text{平方分米}&=(5+0.03)\text{平方米}=5.03\text{平方米}\end{aligned}\]长度与面积的进率不同，不能混用。各原空后已有单位，只填相应数值。
\end{BookAnswer}

% p049_s05
\BookQuestion{一、5}
把0.504、0.45、0.045、0.54、0.405、0.054按从小到大的顺序排列是（\hspace{1.8em}）。\par
\begin{BookAnswer}
\(0.045<0.054<0.405<0.45<0.504<0.54\)。比较时统一写成三位小数：0.045、0.054、0.405、0.450、0.504、0.540，从十分位、百分位、千分位依次比较。
\end{BookAnswer}

% p049_s06
\BookQuestion{一、6}
元元在读一个小数时，把小数点漏读了，结果读成八万四千零二。若原来的小数要读一个零，则原来的小数是（\hspace{1.8em}）。\par
\begin{BookAnswer}
填\(840.02\)。漏读小数点后得到的数字串为84002。若写成8.4002或84.002，小数部分要读两个“零”；8400.2不读“零”。840.02读“八百四十点零二”，恰好只读一个“零”，符合条件。
\end{BookAnswer}

% p049_s07
\BookQuestion{一、7}
在圆圈里填上“＞”“＜”或“＝”。\(6.02\ \bigcirc\ 6.002\)，\allowbreak{}\(3.5\ \bigcirc\ 3.50\)，9千克50克\(\bigcirc\)9.5千克；\allowbreak{}\(7.91\ \bigcirc\ 7.97\)，0.57 m\(\bigcirc\)57 cm，\allowbreak{}\(8.26\ \bigcirc\ 4.62\)。\par
\begin{BookAnswer}
依次填\(>,=,<,<,=,>\)。\par \(6.020>6.002\)，比较百分位；\allowbreak{}\(3.5=3.50\)，小数末尾添0大小不变。不同单位先换算：\[\begin{aligned}50\text{克}&=(50\div1000)\text{千克}=0.05\text{千克}\\9\text{千克}50\text{克}&=(9+0.05)\text{千克}=9.05\text{千克}\\0.57\text{米}&=(0.57\times100)\text{厘米}=57\text{厘米}\end{aligned}\]故9千克50克小于9.5千克，0.57米等于57厘米。其余\(7.91<7.97\)、\(8.26>4.62\)。
\end{BookAnswer}

\Needspace{8\baselineskip}

\pdfbookmark[2]{原书第50页}{orig-50}

\subsection*{原书第50页}

% p050_s01
\BookQuestion{一、8}
一个三位小数“四舍五入”后得到的近似数是7.20，这个三位小数最大是（\hspace{1.8em}），最小是（\hspace{1.8em}）。\par
\begin{BookAnswer}
依次填\(7.204,7.195\)。保留两位小数，要看千分位。\par 从进位想：7.19后面添千分位，最小添5，7.195才会向百分位进1，成为7.20，所以最小7.195。\par 从舍去想：7.20后面添千分位，最大添4，7.204舍去后仍是7.20；若添5就会变7.21，所以最大7.204。\par 在数线上，就是从7.195起、到7.205之前的三位小数。7.205本身不能取。判断依据是规定保留的数位，不是把小数末尾0一律排除。
\end{BookAnswer}

% p050_s02
\BookQuestion{二、1}
下图中不能表示0.06的是（\hspace{1.8em}）。\par
\noindent A. 图A\par
\noindent B. 图B\par
\noindent C. 图C\par
\noindent D. 图D\par
\BookFigure{figures/p050_s02-01.tex}
\begin{BookAnswer}
选\(D\)。A在百分位拨6颗珠，表示6个0.01，即0.06；B的百格中涂6格，表示0.06；C把0到0.1的长度分成5个相等间隔，每格0.02，实心点在0.06处。D把一个整体平均分成10份，涂6份，表示0.6，不是0.06。
\end{BookAnswer}

% p050_s03
\BookQuestion{二、2}
下面的数去掉0以后，大小不变的是（\hspace{1.8em}）。\par
\noindent A. 30.025\par
\noindent B. 32.500\par
\noindent C. 30025\par
\noindent D. 300.25\par
\begin{BookAnswer}
选B。只有小数末尾的0可以去掉而大小不变，\allowbreak{}\(32.500=32.5\)。A、C、D的0都在整数部分或小数中间，删去会改变数位，数的大小随之改变。
\end{BookAnswer}

% p050_s04
\BookQuestion{二、3}
下面的小数中，与9最接近的是（\hspace{1.8em}）。\par
\noindent A. 8.9\par
\noindent B. 9.01\par
\noindent C. 9.001\par
\noindent D. 8.89\par
\begin{BookAnswer}
选C。四个数与9的距离分别是\(0.1,0.01,0.001,0.11\)，其中0.001最小，所以9.001最接近9。
\end{BookAnswer}

% p050_s05
\BookQuestion{二、4}
下图是某小学运动会3名选手的跳远情况，方框里数是（\hspace{1.8em}）。\par
\noindent A. 2\par
\noindent B. 3\par
\noindent C. 4\par
\noindent D. 5\par
\BookFigure{figures/p050_s05-01.tex}
\begin{BookAnswer}
选C，方框中是4，跳远距离为1.43米。图中中间落点比1.24米远、比1.53米近，且更靠近1.53米。若为1.33米，它到1.24的差0.09小于到1.53的差0.20，位置应偏左；若为1.43米，它到两边的差为0.19、0.10，与图上偏右的位置相符。仅凭“在两数之间”不足以区分3和4，还要看图中的相对距离。
\BookFigure{figures/p050_s05-02.tex}
\end{BookAnswer}

% p050_s06
\BookQuestion{二、5}
\(0.\square8<0.6\)，方框里可以填的数有（\hspace{1.8em}）个。\par
\noindent A. 4\par
\noindent B. 5\par
\noindent C. 6\par
\noindent D. 7\par
\begin{BookAnswer}
选C，共6个。\(0.\square8<0.60\)，十分位可取0、1、2、3、4、5；取6时，0.68大于0.60，因此不能取6、7、8、9。
\end{BookAnswer}

% p050_s07
\BookQuestion{二、6}
用6、9、2、0这四个数字和小数点组成不同的三位小数，精确到个位是3的小数有（\hspace{1.8em}）个。\par
\noindent A. 6\par
\noindent B. 5\par
\noindent C. 4\par
\noindent D. 3\par
\begin{BookAnswer}
选C，共4个。精确到个位为3，数应在\(2.5\le x<3.5\)中。四张数字卡各用一次、组成三位小数，只能以2作整数部分，十分位须为6或9。列出\(2.609,2.690,2.906,2.960\)，共4个。末尾0保留的2.690、2.960也都是符合要求的三位小数写法。
\end{BookAnswer}

% p050_s08
\BookQuestion{二、7}
一个小数先扩大到它的100倍，再把小数点向左移动一位后是39.25，这个小数原来是（\hspace{1.8em}）。\par
\noindent A. 392.5\par
\noindent B. 39.25\par
\noindent C. 3.925\par
\noindent D. 0.3925\par
\begin{BookAnswer}
选C。先扩大到100倍，再缩小到十分之一，合起来扩大到\(100\div10=10\)倍。反向求原数：\(39.25\div10=3.925\)。也可先把39.25扩大10倍得392.5，再除以100，仍得3.925。
\end{BookAnswer}

\Needspace{8\baselineskip}

\pdfbookmark[2]{原书第51页}{orig-51}

\subsection*{原书第51页}

% p051_s01
\BookQuestion{三、1}
直接写得数。\(10.56\times10=\)，\allowbreak{}\(3.15\times1000=\)，\allowbreak{}\(2.71\div10=\)；\allowbreak{}\(12.1\div100=\)，\allowbreak{}\(7.69\times100\div10=\)，\allowbreak{}\(10.02\div100\times1000=\)。\par
\begin{BookAnswer}
按题面顺序，得\(105.6,3150,0.271;0.121,76.9,100.2\)。\par 乘10、100、1000时，小数点分别右移1、2、3位；除以这些数时分别左移。\(7.69\times100\div10\)相当于乘10，得76.9；\allowbreak{}\(10.02\div100\times1000\)也相当于乘10，得100.2。
\end{BookAnswer}

% p051_s02
\BookQuestion{三、2}
在圆圈里填上“×”或“÷”，在方框里填上适当的数。\(3.2\bigcirc\square=0.32\)，\allowbreak{}\(2.4\bigcirc\square=240\)，\allowbreak{}\(\square\div10=3.9\)；\allowbreak{}\(0.8\bigcirc\square=8\)，\allowbreak{}\(\square\bigcirc10=6.7\)，\allowbreak{}\(\square\div1000=0.735\)。\par
\begin{BookAnswer}
按本课“小数点移动”的方法，从左到右、先上后下可填：\(3.2\div10=0.32\)，\allowbreak{}\(2.4\times100=240\)，\allowbreak{}\(39\div10=3.9\)；\allowbreak{}\(0.8\times10=8\)，\allowbreak{}\(67\div10=6.7\)，\allowbreak{}\(735\div1000=0.735\)。先比较变化前后小数点移了几位，再确定乘或除10、100、1000。第五式也可填\(0.67\times10=6.7\)，同样属于已学的小数点移动方法。题面未限定方框为整数，故此填写也合理。
\end{BookAnswer}

% p051_s03
\BookQuestion{三、3}
在每道题的方框里填上相邻的两个整数。\(\square<2.5<\square\)，\allowbreak{}\(\square<10.08<\square\)，\allowbreak{}\(\square>6.94>\square\)；\allowbreak{}\(\square>35.618>\square\)，\allowbreak{}\(\square<99.2<\square\)，\allowbreak{}\(\square<275.2<\square\)。\par
\begin{BookAnswer}
依次填\(2,3;10,11;7,6;36,35;99,100;275,276\)。\par 一个非整数小数总位于相邻的两个整数之间：\(2<2.5<3\)，\allowbreak{}\(10<10.08<11\)，\allowbreak{}\(7>6.94>6\)，\allowbreak{}\(36>35.618>35\)，\allowbreak{}\(99<99.2<100\)，\allowbreak{}\(275<275.2<276\)。不等号方向不同，填入的左右整数顺序也要随之变化。
\end{BookAnswer}

% p051_s04
\BookQuestion{四、1}
涂色表示下面各数：0.7、0.39、1.43。（3分）\par
\BookFigure{figures/p051_s04-01.tex}
\begin{BookAnswer}
0.7：圆平均分成10份，涂其中7份。\par 0.39：百格平均分成100份，涂39格，例如3整列加第4列的9格。\par 1.43：先把一个整方形全部涂满，再在后面的百格中涂43格，例如4整列加第5列的3格。前面的整方形与后面的百格，都分别表示同样大小的“1”。
\BookFigure{figures/p051_s04-02.tex}
\end{BookAnswer}

% p051_s05
\BookQuestion{四、2}
在直线上用实心点标出下面各数的位置，并写出对应的数。4.4、4.7、5.0、5.4、5.65。（5分）\par
\BookFigure{figures/p051_s05-01.tex}
\begin{BookAnswer}
先认每一小格表示多少：相邻整数间分成10格，每格0.1。4.4在4右边第4格，4.7在4右边第7格，5.0就是5，5.4在5右边第4格。\par 再找5.65：它在5.6和5.7之间。把这一格再平均分成10份，每份0.01，5.65在5.6右边第5小份，也就是正中间。先定位、再在数轴线上点实心点，数字标注在旁边；5.65不能挪到5.6或5.7刻度上。
\BookFigure{figures/p051_s05-02.tex}
\end{BookAnswer}

% p051_s06
\BookQuestion{五、1（1）}
西瓜中含有丰富的维生素和微量元素，是一种营养丰富的水果。张阿姨在超市买了一个10千克的西瓜。（1）已知张阿姨买这个西瓜花了48元，则这个西瓜每千克多少钱？（5分）\par
\begin{BookAnswer}
总价除以质量，得到每千克价格。\[\begin{aligned}&48\div10\\&=4.8\text{(元)}\end{aligned}\]\par 答：每千克4.8元。\par 核验：\(4.8\times10=48\)元，与已知总价相同。
\end{BookAnswer}

\Needspace{8\baselineskip}

\pdfbookmark[2]{原书第52页}{orig-52}

\subsection*{原书第52页}

% p052_s01
\BookQuestion{五、1（2）}
张阿姨在超市买了一个10千克的西瓜。（2）如果每100克瓜可以榨出81.5克西瓜汁，那么张阿姨买的这个西瓜可以榨出多少千克西瓜汁？（西瓜皮厚度忽略不计）（5分）\par
\begin{BookAnswer}
先统一质量单位：\[10\text{千克}=10000\text{克}\]每100克西瓜可榨81.5克汁。\par 方法一：先求100克的份数，再用每份果汁质量乘份数。\[\begin{aligned}&81.5\times(10000\div100)\\={}&81.5\times100\\={}&8150\text{(克)}\end{aligned}\]\par 方法二：先求1千克西瓜的出汁量，再乘10。\[1\text{千克}=1000\text{克}\]\[\begin{aligned}&81.5\times(1000\div100)\times10\\={}&81.5\times10\times10\\={}&81.5\times(10\times10)\\={}&81.5\times100\\={}&8150\text{(克)}\end{aligned}\]\[8150\text{克}=(8150\div1000)\text{千克}=8.15\text{千克}\]\par 答：可以榨出8.15千克西瓜汁。
\end{BookAnswer}

% p052_s02
\BookQuestion{五、2}
张老师要带东东、明明和亮亮三名同学坐动车去外地参加机器人比赛，成人票每张132元。东东身高135 cm，明明身高119 cm，亮亮身高158 cm。张老师和同学们买单程车票共需要多少元？（6分）\par
\par\smallskip\Needspace{4\baselineskip}\begingroup\small\setlength{\tabcolsep}{3pt}
\noindent\textbf{购票须知}\par\smallskip
\noindent\begin{tabularx}{\linewidth}{|*{2}{>{\centering\arraybackslash}X|}}\hline
身高 & 购票标准\\\hline
1.5 m以上 & 成人票\\\hline
1.20～1.50 m & 半票\\\hline
1.20 m以下 & 免票\\\hline
\end{tabularx}\par
\noindent \textcolor{answerred}{\kaishu 按原题所列标准作答；本题三名同学均不处于边界身高。}\par
\endgroup\smallskip
\begin{BookAnswer}
先统一身高单位：\[\begin{aligned}135\text{厘米}&=1.35\text{米}\\119\text{厘米}&=1.19\text{米}\\158\text{厘米}&=1.58\text{米}\end{aligned}\]厘米化米都除以100。东东1.35米买半票，明明1.19米免票，亮亮1.58米与张老师各买一张成人票。\par 方法一：2张成人票加1张半票。\[\begin{aligned}&132\times2+132\div2\\={}&264+66\\={}&330\text{(元)}\end{aligned}\]答：四人买单程票共需要330元。\par 方法二：半张成人票的钱看作1份，2张成人票相当于4份，再加东东1份，共5份。\[\begin{aligned}&(132\div2)\times(2\times2+1)\\={}&66\times(4+1)\\={}&66\times5\\={}&330\text{(元)}\end{aligned}\]答：四人买单程票共需要330元。
\end{BookAnswer}

% p052_s03
\BookQuestion{五、3}
某小学长方形文化墙每行有18块瓷砖，每列有5块瓷砖。每块瓷砖长40 cm，宽30 cm，文化墙的面积是多少平方米？（6分）\par
\begin{BookAnswer}
方法一：先求每块瓷砖面积，再求全部瓷砖的总面积。每行18块，共5行。\[\begin{aligned}&(40\times30)\times(18\times5)\\&=1200\times90\\&=108000\text{(平方厘米)}\end{aligned}\]\par 方法二：先求文化墙的长和宽，再求面积。\[\begin{aligned}&(40\times18)\times(30\times5)\\&=720\times150\\&=108000\text{(平方厘米)}\end{aligned}\]\par 两种方法所得面积先以平方厘米计算，最后另列单位换算：\[\begin{aligned}1\text{平方米}&=10000\text{平方厘米}\\108000\text{平方厘米}&=(108000\div10000)\text{平方米}\\&=10.8\text{平方米}\end{aligned}\]\par 答：文化墙的面积是10.8平方米。主解先进行整数面积计算，再利用小数点移动换算，不需要使用一般小数乘法。
\end{BookAnswer}

% p052_s04
\BookQuestion{五、4}
用数字卡片1、3、6和小数点，能组成多少个不同的一位小数？将这些小数按照从大到小的顺序排列起来。（6分）\par
\begin{BookAnswer}
按题意三张数字卡各用一次。按十位是1、3、6分成3类，每类都有2个不同的一位小数，所以按“每类个数×类别数”列式：\[\begin{aligned}&2\times3\\={}&6\text{(个)}\end{aligned}\]\par 1开头有13.6、16.3；3开头有31.6、36.1；6开头有61.3、63.1。每一类都列全，三类之间也不会重复。\par 从大到小：\(63.1>61.3>36.1>31.6>16.3>13.6\)。\par 答：能组成6个不同的一位小数，顺序如上。
\end{BookAnswer}

\Needspace{12\baselineskip}

\section{第六单元素养评价（B）}

\Needspace{8\baselineskip}

\pdfbookmark[2]{原书第53页}{orig-53}

\subsection*{原书第53页}

% p053_s01
\BookQuestion{一、1}
和1相邻的一位小数分别是（\hspace{1.8em}）和（\hspace{1.8em}）。\par
\begin{BookAnswer}
依次填\(0.9,1.1\)。一位小数的计数单位是0.1，1可写成1.0；它前后相邻的一位小数分别为\(1.0-0.1=0.9\)、\(1.0+0.1=1.1\)。
\end{BookAnswer}

% p053_s02
\BookQuestion{一、2}
科学家们发现，当外界环境温度达到人体正常温度与0.618的乘积时，人会感到最舒服。0.618的计数单位是（\hspace{1.8em}），它含有（\hspace{1.8em}）个这样的计数单位，如果改写成以万分之一为计数单位的数，这个数是（\hspace{1.8em}）。\par
\begin{BookAnswer}
依次填“千分之一、618、0.6180”。0.618写到千分位，计数单位是\(0.001\)，含有618个这样的单位。\par 最后一问要求改写成以“万分之一”为计数单位，须按小数性质在末尾添一个0，写成四位小数\(0.6180\)。它有6180个\(0.0001\)。\(0.618=0.6180\)，数值没有改变，但按题目指定数位表达时，末尾0用来明确写到了万分位，不能省略。
\end{BookAnswer}

% p053_s03
\BookQuestion{一、3}
图中两个“4”表示的意义不同，这两个“4”相差（\hspace{1.8em}）。\par
\BookFigure{figures/p053_s03-01.tex}
\begin{BookAnswer}
填\(0.36\)。计数器百分位4颗珠表示0.04；圆平均分成10份，涂4份表示0.4。两者相差\(0.4-0.04=0.36\)。同一个数字在不同数位，表示的大小不同。
\end{BookAnswer}

% p053_s04
\BookQuestion{一、4}
在0.03、0.003、0.030这三个数中，计数单位相同的是（\hspace{1.8em}）和（\hspace{1.8em}）；大小相等的是（\hspace{1.8em}）和（\hspace{1.8em}）。\par
\begin{BookAnswer}
计数单位相同的是\(0.003\)与\(0.030\)，大小相等的是\(0.03\)与\(0.030\)。前两者写成三位小数，计数单位均为千分之一；0.03写成两位小数，计数单位为百分之一。小数末尾添0大小不变，所以0.03=0.030。必须区分“数的大小”和“指定写法的计数单位”。
\end{BookAnswer}

% p053_s05
\BookQuestion{一、5}
一个数由3个十，8个十分之一和5个千分之一组成，这个数写作（\hspace{1.8em}），读这个数时，有（\hspace{1.8em}）个零不必读出来。\par
\begin{BookAnswer}
依次填\(30.805,1\)。3个十是30，8个十分之一是0.8，5个千分之一是0.005，合成30.805。整数部分30末尾的0不读，小数部分的0要逐位读，读作“三十点八零五”，所以有1个零不读。
\end{BookAnswer}

% p053_s06
\BookQuestion{一、6}
甲数是3.8，乙数是38，在它们的末尾都填上两个零，这时乙数是甲数的（\hspace{1.8em}）倍。\par
\begin{BookAnswer}
填\(1000\)。甲的小数末尾添两个0，3.8变成3.800，大小不变；乙的整数末尾添两个0，38变成3800，扩大100倍。乙原来是甲的10倍，变化后是\(10\times100=1000\)倍。
\end{BookAnswer}

% p053_s07
\BookQuestion{一、7}
一个三位小数，各个数位上的数字之和是9，且各个数位上的数字都不相同，符合条件的三位小数最小是（\hspace{1.8em}）。\par
\begin{BookAnswer}
填\(0.126\)。“各个数位上的数字都不相同”包括个位。先让个位为0，十分位最小只能取1，百分位最小取2，千分位由数字和决定：\(9-0-1-2=6\)。所以最小为0.126。0.018虽然更小，但个位与十分位都为0，不符合数位数字互不相同的条件。
\end{BookAnswer}

% p053_s08
\BookQuestion{一、8}
在括号里填上合适的单位，使等式成立。45（\hspace{1.8em}）＝0.45（\hspace{1.8em}）；3540（\hspace{1.8em}）＝3.540（\hspace{1.8em}）。\par
\begin{BookAnswer}
一组答案为“厘米、米；克、千克”。前式进率100，后式进率1000：\[\begin{aligned}45\text{厘米}&=(45\div100)\text{米}=0.45\text{米}\\3540\text{克}&=(3540\div1000)\text{千克}=3.540\text{千克}\end{aligned}\]题目未限定单位种类，原有其他合理组合也分别列出：\[\begin{aligned}45\text{分}&=0.45\text{元}\\45\text{平方厘米}&=0.45\text{平方分米}\\45\text{平方分米}&=0.45\text{平方米}\\45\text{公顷}&=0.45\text{平方千米}\\3540\text{毫米}&=3.540\text{米}\\3540\text{米}&=3.540\text{千米}\\3540\text{千克}&=3.540\text{吨}\end{aligned}\]保留题面给定的3.540写法。两边必须是同一种量，不能把长度与质量混配。这里只在原单位空填写单位，不重复数值。
\end{BookAnswer}

% p053_s09
\BookQuestion{一、9}
10千克花生可以榨4.5千克花生油，1吨花生可以榨（\hspace{1.8em}）千克花生油，合（\hspace{1.8em}）吨花生油。（第二空接下页）\par
\begin{BookAnswer}
依次填\(450,0.45\)。\[1\text{吨}=1000\text{千克}\]是10千克的100倍。\[\begin{aligned}&4.5\times(1000\div10)\\&=4.5\times100\\&=450\text{(千克)}\end{aligned}\]\[450\text{千克}=(450\div1000)\text{吨}=0.45\text{吨}\]
\end{BookAnswer}

\Needspace{8\baselineskip}

\pdfbookmark[2]{原书第54页}{orig-54}

\subsection*{原书第54页}

% p054_s01
\BookQuestion{一、10}
若\(A\)表示1，\allowbreak{}\(B\)表示2，则\(C\)表示（\hspace{1.8em}）；若\(A\)表示2.3，\allowbreak{}\(B\)表示2.4，则\(D\)表示（\hspace{1.8em}）；若\(A\)表示48.6，\allowbreak{}\(B\)表示48.7，则\(C\)表示（\hspace{1.8em}）。\par
\BookFigure{figures/p054_s01-01.tex}
\begin{BookAnswer}
依次填\(1.7,2.42,48.67\)。\(A\)到\(B\)有10小格，\allowbreak{}\(A\)到\(C\)有7格，\allowbreak{}\(B\)到\(D\)有2格。\par 当\(A=1,B=2\)时，1平均分成10份，每格表示0.1；\allowbreak{}\(C\)比\(A\)多7个0.1，所以是1.7。\par 当\(A=2.3,B=2.4\)时，0.1平均分成10份，每格表示0.01；\allowbreak{}\(D\)比\(B\)多2个0.01，所以是2.42。\par 当\(A=48.6,B=48.7\)时，每格表示0.01；\allowbreak{}\(C\)比\(A\)多7个0.01，所以是48.67。
\end{BookAnswer}

% p054_s02
\BookQuestion{二、1}
三星堆出土的青铜神树高3.96米，是全球最大单件青铜文物。3.96米中的“9”表示（\hspace{1.8em}）。\par
\noindent A. 9米\par
\noindent B. 0.9分米\par
\noindent C. 9分米\par
\noindent D. 0.9厘米\par
\begin{BookAnswer}
选C。9在十分位，表示0.9米。米化分米的进率是10：\[0.9\text{米}=(0.9\times10)\text{分米}=9\text{分米}\]不能把9理解成9米或0.9分米。
\end{BookAnswer}

% p054_s03
\BookQuestion{二、2}
下列说法正确的是（\hspace{1.8em}）。\par
\noindent A. 用“四舍五入”法保留两位小数时，关键要看百分位上的数\par
\noindent B. 一个两位小数精确到十分位约是3.5，这个数最小是3.50\par
\noindent C. 一个小数的小数位数越多，这个小数就越大\par
\noindent D. 所有的小数都有整数部分\par
\begin{BookAnswer}
选D。所有小数都有整数部分，整数部分可以是0。A保留两位小数应看千分位，而不是百分位；B最小可为3.45，四舍五入到十分位也是3.5，并非最小3.50；C数位越多不一定越大，例如\(0.99<1.0\)。
\end{BookAnswer}

% p054_s04
\BookQuestion{二、3}
计算\(0.2\times0.03=(2\times0.1)\times(3\times0.01)=(2\times3)\times(0.1\times0.01)=6\times0.001=0.006\)，第三步的“6”表示的是新的（\hspace{1.8em}），“0.001”表示的是新的（\hspace{1.8em}）。\par
\noindent A. 计数单位\par
\noindent B. 计数单位的个数\par
\noindent C. 积的小数位数\par
\noindent D. 积\par
\begin{BookAnswer}
依次选\(\mathrm{B},\mathrm{A}\)。由题中给出的第三步\(6\times0.001=0.006\)，可知积是\(6\)个\(0.001\)；所以“6”表示新的计数单位的个数，“0.001”表示新的计数单位，即千分之一。\par 本题只解读已经给出的算式，不要求学生在本阶段掌握一般小数乘法的计算方法。
\end{BookAnswer}

% p054_s05
\BookQuestion{二、4}
学校运动会上，三名同学参加100米短跑比赛，成绩如下表。根据比赛排名，浩浩获得了一等奖，那么方框里最大能填（\hspace{1.8em}）。\par
\noindent A. 3\par
\noindent B. 4\par
\noindent C. 5\par
\noindent D. 6\par
\par\smallskip\Needspace{4\baselineskip}\begingroup\small\setlength{\tabcolsep}{3pt}
\noindent\begin{tabularx}{\linewidth}{|*{4}{>{\centering\arraybackslash}X|}}\hline
选手 & 明明 & 亮亮 & 浩浩\\\hline
成绩（秒） & 15.63 & 15.47 & \(15.\square6\)\\\hline
\end{tabularx}\par
\endgroup\smallskip
\begin{BookAnswer}
选B，方框中最大填4。跑步用时越短，名次越靠前。浩浩要得第一，\allowbreak{}\(15.\square6<15.47\)。填4时15.46比15.47小；填5时15.56已比15.47大，因此最大填4。
\end{BookAnswer}

% p054_s06
\BookQuestion{二、5}
计算“8米6厘米＋560厘米”时，用“米”作单位，正确的算式是（\hspace{1.8em}）。\par
\noindent A. \(8.6+5.6\)\par
\noindent B. \(8.06+5.6\)\par
\noindent C. \(8.06+5.06\)\par
\noindent D. \(8.6+5.06\)\par
\begin{BookAnswer}
选B。先把长度都化成米：\[\begin{aligned}6\text{厘米}&=(6\div100)\text{米}=0.06\text{米}\\8\text{米}6\text{厘米}&=(8+0.06)\text{米}=8.06\text{米}\\560\text{厘米}&=(560\div100)\text{米}=5.6\text{米}\end{aligned}\]正确算式为\(8.06+5.6\)。6厘米是0.06米，560厘米是5.6米，不能把厘米数直接接在小数点后。
\end{BookAnswer}

% p054_s07
\BookQuestion{三、1}
直接写得数。\(40\div100=\)，\allowbreak{}\(0.6\times1000=\)，\allowbreak{}\(3.6\div10=\)，\allowbreak{}\(45\div10\times100=\)；\allowbreak{}\(0.21\div10=\)，\allowbreak{}\(250\div1000=\)，\allowbreak{}\(4800\div100=\)，\allowbreak{}\(3.2\times1000\div10=\)。\par
\begin{BookAnswer}
按题面顺序，得\(0.4,600,0.36,450;0.021,0.25,48,320\)。\par \(40\div100=0.4\)；\allowbreak{}\(0.6\times1000=600\)；\allowbreak{}\(3.6\div10=0.36\)；\allowbreak{}\(45\div10\times100=45\times(100\div10)=45\times10=450\)。\par \(0.21\div10=0.021\)；\allowbreak{}\(250\div1000=0.25\)；\allowbreak{}\(4800\div100=48\)；\allowbreak{}\(3.2\times1000\div10=3.2\times(1000\div10)=3.2\times100=320\)。连续放大、缩小可把倍数合并，但必须保持乘除的意义。
\end{BookAnswer}

\Needspace{8\baselineskip}

\pdfbookmark[2]{原书第55页}{orig-55}

\subsection*{原书第55页}

% p055_s01
\BookQuestion{三、2（1）}
求下列各数的近似数。精确到百分位：\(2.9633\approx(\quad)\)，\allowbreak{}\(249.3370\approx(\quad)\)。\par
\begin{BookAnswer}
精确到百分位，要看千分位。\(2.9633\)的千分位是3，舍去；\allowbreak{}\(249.3370\)的千分位是7，百分位进1。因此填\(2.96\)、\(249.34\)。四舍五入看原数中保留末位的下一位：小于5舍去，大于或等于5向前进1，不能先保留三位再连续取近似数。
\end{BookAnswer}

% p055_s02
\BookQuestion{三、2（2）}
先改写成用“万”或“亿”作单位的数，再保留一位小数。\(327900=(\quad)\text{万}\approx(\quad)\text{万}\)；\allowbreak{}\(25499270000=(\quad)\text{亿}\approx(\quad)\text{亿}\)。\par
\begin{BookAnswer}
改写成“万”作单位时除以10000，改写成“亿”作单位时除以100000000。\[\begin{aligned}327900&=32.79\text{万}\\32.79\text{万}&\approx32.8\text{万}\\25499270000&=254.9927\text{亿}\\254.9927\text{亿}&\approx255.0\text{亿}\end{aligned}\]保留一位小数时，\allowbreak{}\(255.0\)末尾的0必须保留，表示精确程度。改写单位没有改变数的大小，用“=”；求近似数才用“≈”。保留一位小数须观察百分位，\allowbreak{}\(254.9927\)的百分位9使十分位进1，连续进位后写\(255.0\)，不能写成255。
\end{BookAnswer}

% p055_s03
\BookQuestion{三、3（1）}
\(\square\)里最大能填几？\(3.\square>3.68\)。\par
\begin{BookAnswer}
\(3.\square>3.68\)。只填一位数字时，最大可填9；\allowbreak{}\(3.9>3.68\)。
\end{BookAnswer}

% p055_s04
\BookQuestion{三、3（2）}
\(\square\)里最大能填几？\(2.07>2.\square6\)。\par
\begin{BookAnswer}
\(2.07>2.\square6\)。十分位只能填0；\allowbreak{}\(2.06<2.07\)，若填1便有\(2.16>2.07\)。填0。
\end{BookAnswer}

% p055_s05
\BookQuestion{三、3（3）}
\(\square\)里最大能填几？\(3240\text{克}>\square.15\text{千克}\)。\par
\begin{BookAnswer}
先统一单位，克化千克除以1000：\[3240\text{克}=(3240\div1000)\text{千克}=3.24\text{千克}\]再比较\(3.24>\square.15\)，整数部分最大填3。
\end{BookAnswer}

% p055_s06
\BookQuestion{三、3（4）}
\(\square\)里最大能填几？\(6.25\text{吨}>6\square25\text{千克}\)。\par
\begin{BookAnswer}
吨化千克乘1000：\[6.25\text{吨}=(6.25\times1000)\text{千克}=6250\text{千克}\]再比较\(6250>6\square25\)，百位最大填2，因为\(6225<6250\)，而\(6325>6250\)。
\end{BookAnswer}

% p055_s07
\BookQuestion{三、3（5）}
\(\square\)里最大能填几？\(0.\square5\text{米}>15\text{厘米}\)。\par
\begin{BookAnswer}
厘米化米除以100：\[15\text{厘米}=(15\div100)\text{米}=0.15\text{米}\]再比较\(0.\square5>0.15\)，十分位最大可填9。
\end{BookAnswer}

% p055_s08
\BookQuestion{三、3（6）}
\(\square\)里最大能填几？\(21.5>21.\square9\)。\par
\begin{BookAnswer}
\(21.5=21.50\)。要使\(21.50>21.\square9\)，十分位最大填4；\allowbreak{}\(21.49<21.50\)，而\(21.59>21.50\)。
\end{BookAnswer}

% p055_s09
\BookQuestion{四、1}
涂色表示下面的小数：\(1.4\)，\allowbreak{}\(1.07\)。每个完整正方形表示1。\par
\BookFigure{figures/p055_s09-01.tex}
\begin{BookAnswer}
\(1.4=1+0.4\)：先把一个完整正方形涂满，再把另一个正方形平均分成10份，涂其中4份。\par \(1.07=1+0.07\)：涂满一个完整正方形，再在百格图中涂7小格；十分位是0，不涂十分之一条。涂色端点应落在原格线上。
\BookFigure{figures/p055_s09-02.tex}
\end{BookAnswer}

% p055_s10
\BookQuestion{四、2（1）}
如下表是4名同学的身高信息。根据线段图回答问题。（1）箭头所表示的是（\hspace{1.8em}）的身高。\par
\par\smallskip\Needspace{4\baselineskip}\begingroup\small\setlength{\tabcolsep}{3pt}
\noindent\begin{tabularx}{\linewidth}{|*{5}{>{\centering\arraybackslash}X|}}\hline
学生 & 乐乐 & 王华 & 可可 & 冬冬\\\hline
身高/米 & 1.32 & 1.19 & 1.45 & \(1.\square4\)\\\hline
\end{tabularx}\par
\endgroup\smallskip
\BookFigure{figures/p055_s10-01.tex}
\begin{BookAnswer}
数轴上相邻两大刻度相差0.1，中间有10小格，每小格是0.01。箭头指向1.45，表中可可是1.45米，因此填“可可”。
\end{BookAnswer}

% p055_s11
\BookQuestion{四、2（2）}
如下表是4名同学的身高信息。根据线段图回答问题。（2）在图中标出乐乐和王华的身高。\par
\par\smallskip\Needspace{4\baselineskip}\begingroup\small\setlength{\tabcolsep}{3pt}
\noindent\begin{tabularx}{\linewidth}{|*{5}{>{\centering\arraybackslash}X|}}\hline
学生 & 乐乐 & 王华 & 可可 & 冬冬\\\hline
身高/米 & 1.32 & 1.19 & 1.45 & \(1.\square4\)\\\hline
\end{tabularx}\par
\endgroup\smallskip
\BookFigure{figures/p055_s11-01.tex}
\begin{BookAnswer}
乐乐1.32米，在1.3右边第2小格；王华1.19米，在1.2左边第1小格。用箭头或点准确落在相应小刻度，再标姓名。
\BookFigure{figures/p055_s11-02.tex}
\end{BookAnswer}

% p055_s12
\BookQuestion{四、2（3）}
如下表是4名同学的身高信息。根据线段图回答问题。（3）冬冬的身高从高到低排第三，他的身高是（\hspace{1.8em}）米。\par
\par\smallskip\Needspace{4\baselineskip}\begingroup\small\setlength{\tabcolsep}{3pt}
\noindent\begin{tabularx}{\linewidth}{|*{5}{>{\centering\arraybackslash}X|}}\hline
学生 & 乐乐 & 王华 & 可可 & 冬冬\\\hline
身高/米 & 1.32 & 1.19 & 1.45 & \(1.\square4\)\\\hline
\end{tabularx}\par
\endgroup\smallskip
\BookFigure{figures/p055_s12-01.tex}
\begin{BookAnswer}
冬冬身高为\(1.\square4\)米。按从高到低第三名，它必须低于乐乐1.32米而高于王华1.19米，所以十分位只能是2，身高为1.24米。次序：\(1.45>1.32>1.24>1.19\)。
\end{BookAnswer}

% p055_s13
\BookQuestion{五、1}
一块正方形菜地的边长缩小到原来的\(\frac1{100}\)之后是\(0.03\)米，原来这块菜地的面积是多少平方米？\par
\begin{BookAnswer}
原边长是缩小后边长的100倍。\[\begin{aligned}&0.03\times100\\={}&3\text{(米)}\end{aligned}\]再求原面积：\[\begin{aligned}&3\times3\\={}&9\text{(平方米)}\end{aligned}\]答：原来这块菜地的面积是9平方米。
\end{BookAnswer}

\Needspace{8\baselineskip}

\pdfbookmark[2]{原书第56页}{orig-56}

\subsection*{原书第56页}

% p056_s01
\BookQuestion{五、2}
快递公司把\(3.75\)吨大件快递平均分给5辆货车运送，每辆货车装20件，平均每件快递重多少千克？像这样的1万件快递重多少千克？\par
\begin{BookAnswer}
先统一单位，吨化千克乘1000：\[3.75\text{吨}=(3.75\times1000)\text{千克}=3750\text{千克}\]方法一：先求5辆车共装多少件，再求每件质量。\[\begin{aligned}&3750\div(20\times5)\\={}&3750\div100\\={}&37.5\text{(千克)}\end{aligned}\]方法二：先分车再分件的数量分析仍然成立；计算时用连续除法的性质，把5与20先凑成100，直接除以100。\[\begin{aligned}&3750\div5\div20\\={}&3750\div(5\times20)\\={}&3750\div100\\={}&37.5\text{(千克)}\end{aligned}\]答：平均每件快递重37.5千克。\par 先改写件数：\[1\text{万件}=10000\text{件}\]1万件的质量：\[\begin{aligned}&37.5\times10000\\={}&375000\text{(千克)}\end{aligned}\]答：像这样的1万件快递重375000千克。可检验：每件\(37.5\)千克，\allowbreak{}\(20\times5=100\)件，共\(37.5\times100=3750\)千克，与题目总质量相同。
\end{BookAnswer}

% p056_s02
\BookQuestion{五、3（1）}
单板滑雪男子决赛，下表是前三名选手比赛得分。（1）若根据三轮比赛最高得分排名，那么（\hspace{1.8em}）号荣获第一名。\par
\par\smallskip\Needspace{4\baselineskip}\begingroup\small\setlength{\tabcolsep}{3pt}
\noindent\begin{tabularx}{\linewidth}{|*{4}{>{\centering\arraybackslash}X|}}\hline
选手 & 第一轮 & 第二轮 & 第三轮\\\hline
1号 & 28.75 & 94.00 & 76.25\\\hline
2号 & 89.00 & 33.00 & 93.5\\\hline
3号 & 89.25 & 75.75 & 82.00\\\hline
\end{tabularx}\par
\endgroup\smallskip
\begin{BookAnswer}
三名选手单轮最高分分别为94分、93.5分、89.25分。\(94>93.5>89.25\)，按三轮中的最高得分排名，1号第一。
\end{BookAnswer}

% p056_s03
\BookQuestion{五、3（2）}
单板滑雪男子决赛，下表是前三名选手比赛得分。（2）若根据三轮比赛总得分排名，那么（\hspace{1.8em}）号荣获第一名。\par
\par\smallskip\Needspace{4\baselineskip}\begingroup\small\setlength{\tabcolsep}{3pt}
\noindent\begin{tabularx}{\linewidth}{|*{4}{>{\centering\arraybackslash}X|}}\hline
选手 & 第一轮 & 第二轮 & 第三轮\\\hline
1号 & 28.75 & 94.00 & 76.25\\\hline
2号 & 89.00 & 33.00 & 93.5\\\hline
3号 & 89.25 & 75.75 & 82.00\\\hline
\end{tabularx}\par
\endgroup\smallskip
\begin{BookAnswer}
分别求三轮总分。1号：\[\begin{aligned}&28.75+94+76.25\\={}&(28.75+76.25)+94\\={}&105+94\\={}&199\text{(分)}\end{aligned}\]2号：\[\begin{aligned}&89+33+93.5\\={}&122+93.5\\={}&215.5\text{(分)}\end{aligned}\]3号：\[\begin{aligned}&89.25+75.75+82\\={}&165+82\\={}&247\text{(分)}\end{aligned}\]\(247>215.5>199\)。答：按三轮总分排名，3号第一。
\end{BookAnswer}

% p056_s04
\BookQuestion{五、3（3）}
单板滑雪男子决赛，下表是前三名选手比赛得分。（3）若本次比赛获得冠军的是2号，那么排名的规则是什么？\par
\par\smallskip\Needspace{4\baselineskip}\begingroup\small\setlength{\tabcolsep}{3pt}
\noindent\begin{tabularx}{\linewidth}{|*{4}{>{\centering\arraybackslash}X|}}\hline
选手 & 第一轮 & 第二轮 & 第三轮\\\hline
1号 & 28.75 & 94.00 & 76.25\\\hline
2号 & 89.00 & 33.00 & 93.5\\\hline
3号 & 89.25 & 75.75 & 82.00\\\hline
\end{tabularx}\par
\endgroup\smallskip
\begin{BookAnswer}
一种合理规则：去掉每位选手的一轮最低分，再比较其余两轮总分。1号：\[\begin{aligned}&94+76.25\\={}&170.25\text{(分)}\end{aligned}\]2号：\[\begin{aligned}&89+93.5\\={}&182.5\text{(分)}\end{aligned}\]3号：\[\begin{aligned}&89.25+82\\={}&171.25\text{(分)}\end{aligned}\]\(182.5>171.25>170.25\)。答：按这条规则，2号获得冠军。题目是开放题，给出规则时必须对三名选手使用同一标准。
\end{BookAnswer}

% p056_s05
\BookQuestion{五、4}
A、B、C、D四名同学测量身高，测得的结果是\(138.5\)厘米、\(142.6\)厘米、141厘米、\(138.8\)厘米。已知A比D高，但比B矮，C比D矮。你知道他们的身高各是多少吗？\par
\begin{BookAnswer}
把文字关系排成一列：A比D高，A比B矮，C比D矮，因此\(B>A>D>C\)。四个数从大到小为\(142.6>141>138.8>138.5\)。逐一对应得：B为142.6厘米，A为141厘米，D为138.8厘米，C为138.5厘米。\par 答：A、B、C、D分别为141厘米、142.6厘米、138.5厘米、138.8厘米。也可先确定最高是B、最矮是C，再比较中间的A和D。
\end{BookAnswer}
